\documentclass[11pt,a4paper]{article}

\usepackage[margin=1in]{geometry}
\usepackage[T1]{fontenc}
\usepackage[utf8]{inputenc}
\usepackage{lmodern}
\usepackage{amsmath,amssymb,bm,mathtools}
\usepackage{booktabs}
\usepackage{graphicx}
\usepackage{longtable}
\usepackage{enumitem}
\usepackage{xcolor}
\usepackage{hyperref}
\usepackage{url}

\hypersetup{
  colorlinks=true,
  linkcolor=blue!50!black,
  citecolor=blue!50!black,
  urlcolor=blue!50!black
}

\setlength{\parindent}{0pt}
\setlength{\parskip}{0.45em}
\setlist[itemize]{leftmargin=1.4em,itemsep=0.15em,topsep=0.25em}
\setlist[enumerate]{leftmargin=1.6em,itemsep=0.15em,topsep=0.25em}

\newcommand{\dd}{\mathrm{d}}
\newcommand{\ii}{\mathrm{i}}
\newcommand{\ee}{\mathrm{e}}
\newcommand{\vpar}{v_{\parallel}}
\newcommand{\kperp}{k_{\perp}}
\newcommand{\bvec}{\bm b}
\newcommand{\Bvec}{\bm B}
\newcommand{\Evec}{\bm E}
\newcommand{\grad}{\nabla}
\newcommand{\avg}[1]{\left\langle #1 \right\rangle}
\newcommand{\fluxavg}[1]{\left\langle #1 \right\rangle_{\psi}}
\newcommand{\RHS}{\mathcal{R}}
\newcommand{\Dop}{\mathcal{D}}
\newcommand{\Jzero}{J_0}
\newcommand{\Gzero}{\Gamma_0}
\newcommand{\Lref}{R_{\mathrm{ref}}}
\newcommand{\vref}{v_{\mathrm{th,ref}}}
\newcommand{\rhoref}{\rho_{\mathrm{ref}}}
\newcommand{\crefpath}[1]{\path{#1}}
\DeclareMathOperator{\diag}{diag}
\DeclareMathOperator{\argmax}{arg\,max}

\title{Physical Model and Numerical Scheme for a Differentiable Flux-Tube Stellarator Gyrokinetic Solver}
\author{Mohsen Sadr}
\date{\today}

\begin{document}
\maketitle

\begin{abstract}
This document specifies the first implementation target for the project: a differentiable, local, flux-tube, linear electrostatic gyrokinetic solver for stellarator geometry.  The physics, normalization, signs, and benchmarks follow the published GKW model, the original GKW Fortran source, and the differentiable Gyaradax implementation.  The numerical target follows \path{task.tex}: Fourier discretization in the perpendicular directions and spectral discretization along the magnetic field and velocity-space coordinates.  GKW-style finite differences are retained as a fallback and parity backend.  Text is intentionally brief; equations, array contracts, numerical operators, and tests are written in enough detail to guide implementation in \path{src/}.
\end{abstract}

\tableofcontents

\section{Scope and Reference Hierarchy}

The first solver shall implement the collisionless linear electrostatic local gyrokinetic system in a stellarator flux tube.  The baseline electron model is adiabatic electrons with kinetic ions.  Kinetic electrons, nonlinear ExB dynamics, collisions, electromagnetic perturbations, rotation, and global radial effects are documented as extensions, not as the first implementation target.

When conventions are ambiguous, use the following order of authority.
\begin{enumerate}
  \item Published GKW paper: \path{papers/gkw/GKW.pdf}.
  \item Original GKW source: \path{relevant-codes/gkw/src/}.
  \item Searchable paper reconstruction: \path{papers/gkw/GKW_rebuilt.tex}.
  \item GKW sample cases: \path{relevant-codes/gkw/samples/}.
  \item Gyaradax JAX implementation and tests: \path{relevant-codes/gyaradax/}.
  \item DESC and stellarator optimization papers for Boozer geometry and design-loop coupling.
\end{enumerate}

The implementation must expose a pure residual
\begin{equation}
  \partial_t f = \RHS(f;\phi,\mathcal{G},p),
  \qquad
  \phi = \Phi(f;\mathcal{G},p),
  \label{eq:residual-contract}
\end{equation}
where \(f\) is the gyrocenter distribution perturbation, \(\phi\) is the electrostatic potential, \(\mathcal{G}\) is the flux-tube geometry, and \(p\) collects species, profile, grid, and solver parameters.  The residual is the central object for time stepping, eigenvalue solves, and automatic differentiation.

The primary discretization is:
\begin{equation}
  (x,y)\hbox{ Fourier}
  \quad+\quad
  z\hbox{ spectral}
  \quad+\quad
  (\vpar,\mu)\hbox{ spectral}.
  \label{eq:primary-discretization}
\end{equation}
Finite-difference stencils from GKW/Gyaradax are not the default target for new solver kernels.  They are a documented fallback for source-code parity, debugging, and convergence comparisons.

\subsection{GKW Source Crosswalk}

\begin{longtable}{p{0.31\textwidth}p{0.61\textwidth}}
\toprule
Reference file & Use in this project \\
\midrule
\path{linart.f90} & Main program flow and run sequencing. \\
\path{normalise.F90} & Normalization, growth-rate/frequency conventions. \\
\path{geom.F90} & Geometry, metric elements, drift tensors, circular and general geometry. \\
\path{mode.F90} & Fourier mode grids, radial/binormal mode labels, twist-and-shift connectivity. \\
\path{grid.F90}, \path{velocitygrid.F90} & Legacy grid sizes, velocity-space nodes, and weights for parity/fallback tests. \\
\path{components.F90} & Species parameters, Maxwellians, charges, masses, density and temperature ratios. \\
\path{functions.F90}, \path{specfun.f90} & Special functions used by Bessel and polarization factors. \\
\path{linear_terms.F90} & Linear gyrokinetic terms, Poisson/quasineutrality pieces, sign conventions. \\
\path{non_linear_terms.F90} & Future nonlinear ExB bracket, FFT/dealiasing, nonlinear timestep estimate. \\
\path{exp_integration.F90} & Explicit RK schemes, RHS assembly, field update flow. \\
\path{matdat.F90} & Matrix preparation, fallback stencil coefficients, linear timestep estimates. \\
\path{diagnostic.F90} & Growth rates, frequencies, fluxes, spectra, and output normalization. \\
\path{collisionop.F90}, \path{rotation.F90} & Future extensions beyond the first collisionless non-rotating baseline. \\
\bottomrule
\end{longtable}

\section{Normalization}

All implemented equations should use a dimensionless normalization close to GKW/Gyaradax.  Let \(\Lref\), \(B_{\mathrm{ref}}\), \(m_{\mathrm{ref}}\), \(T_{\mathrm{ref}}\), \(n_{\mathrm{ref}}\), and \(Z_{\mathrm{ref}}e\) be reference values.  Define
\begin{equation}
  \vref = \sqrt{\frac{2T_{\mathrm{ref}}}{m_{\mathrm{ref}}}},
  \qquad
  t_{\mathrm{ref}} = \frac{\Lref}{\vref},
  \qquad
  \rhoref = \frac{m_{\mathrm{ref}}\vref}{Z_{\mathrm{ref}} e B_{\mathrm{ref}}}.
\end{equation}
The dimensionless variables are
\begin{equation}
  \hat t = \frac{t}{t_{\mathrm{ref}}},
  \qquad
  \hat{\bm x} = \frac{\bm x}{\Lref},
  \qquad
  \hat B = \frac{B}{B_{\mathrm{ref}}},
  \qquad
  \hat \phi = \frac{Z_{\mathrm{ref}} e \phi}{T_{\mathrm{ref}}},
  \qquad
  \hat f_s = \frac{f_s}{n_{\mathrm{ref}}/\vref^3}.
\end{equation}
The hats are dropped below.  Species parameters are stored as
\begin{equation}
  \hat m_s = \frac{m_s}{m_{\mathrm{ref}}},
  \quad
  \hat Z_s = \frac{Z_s}{Z_{\mathrm{ref}}},
  \quad
  \hat T_s = \frac{T_s}{T_{\mathrm{ref}}},
  \quad
  \hat n_s = \frac{n_s}{n_{\mathrm{ref}}},
  \quad
  v_{\mathrm{th},s} = \sqrt{\frac{\hat T_s}{\hat m_s}}.
\end{equation}
The local profile gradients are dimensionless inverse scale lengths
\begin{equation}
  \omega_{n,s} = \frac{\Lref}{L_{n,s}}
  = -\Lref \frac{\dd \ln n_s}{\dd r},
  \qquad
  \omega_{T,s} = \frac{\Lref}{L_{T,s}}
  = -\Lref \frac{\dd \ln T_s}{\dd r}.
  \label{eq:profile-gradients}
\end{equation}
Implementation names should be close to Gyaradax/GKW: \(\omega_{n,s}\leftrightarrow\) \texttt{rln}, \(\omega_{T,s}\leftrightarrow\) \texttt{rlt}, \(\hat m_s\leftrightarrow\) \texttt{mas}, \(\hat Z_s\leftrightarrow\) \texttt{signz}, \(\hat T_s\leftrightarrow\) \texttt{tmp}, and \(\hat n_s\leftrightarrow\) \texttt{de}.

\section{Coordinates and Phase Space}

\subsection{Flux-tube coordinates}

Use a local field-aligned coordinate system
\begin{equation}
  (x,y,z) \equiv (\hbox{radial},\hbox{binormal},\hbox{parallel}),
  \qquad
  \Bvec = \grad \psi \times \grad \alpha,
  \label{eq:clebsch}
\end{equation}
where \(\psi\) labels the flux surface and \(\alpha\) labels a field line.  In Boozer angles \((\theta,\varphi)\), a standard choice is
\begin{equation}
  \alpha = \theta - \iota(\psi)\varphi - \alpha_0,
  \qquad
  \iota = \frac{1}{q}.
  \label{eq:alpha-boozer}
\end{equation}
The coordinate \(z\) parameterizes distance along the selected field line.  For compatibility with GKW/Gyaradax, a cell-centered grid on \(z\in[-n_{\mathrm{period}}+1/2,n_{\mathrm{period}}-1/2]\) is acceptable for analytic tests.  For stellarators, \(z\) may be Boozer toroidal angle, Boozer poloidal angle, or an arc-length-like coordinate, provided the geometry supplies \(\bvec\cdot\grad z\).

The Fourier representation in the perpendicular plane is
\begin{equation}
  h(x,y,z,\vpar,\mu,t)
  =
  \sum_{k_x,k_y} \hat h_{k_x,k_y}(z,\vpar,\mu,t)
  \exp\!\left[\ii(k_x x+k_y y)\right],
  \label{eq:fourier}
\end{equation}
with real-field Hermitian symmetry represented by a nonnegative \(k_y\) half spectrum.

\subsection{Phase-space state}

The evolved variable for species \(s\) is
\begin{equation}
  f_s = f_s(\vpar,\mu,z,k_x,k_y,t),
  \label{eq:state}
\end{equation}
where \(\mu = v_\perp^2/(2B)\) in physical variables and is represented on a positive grid.  The first implementation may use one kinetic ion species with adiabatic electrons:
\begin{equation}
  f \in \mathbb{C}^{N_{v}\times N_{\mu}\times N_z\times N_x\times N_y}.
\end{equation}
The multi-species extension adds a leading species axis:
\begin{equation}
  f \in \mathbb{C}^{N_s^{\mathrm{sp}}\times N_{v}\times N_{\mu}\times N_z\times N_x\times N_y}.
\end{equation}

\section{Geometry Required by the Solver}

At every \(z_j\), the solver needs a geometry object \(\mathcal{G}(z_j)\) that supplies enough information to form \(k_\perp^2\), parallel streaming, mirror force, ExB response, and magnetic drifts.  The minimal internal geometry contract is
\begin{equation}
  \mathcal{G} =
  \left\{
  B,\;
  F,\;
  G,\;
  E_y,\;
  D_x,D_y,\;
  g^{xx},g^{xy},g^{yy},\;
  w_z,\;
  \hbox{mode connectivity}
  \right\}.
  \label{eq:geometry-contract}
\end{equation}
This maps naturally to Gyaradax/GKW-style arrays:
\begin{center}
\begin{tabular}{ll}
\toprule
This document & Gyaradax/GKW-like name \\
\midrule
\(B(z)\) & \texttt{bn} \\
\(F(z)\), parallel streaming coefficient & \texttt{ffun} \\
\(G(z) = -B^{-1}\partial_z B\) up to coordinate factors & \texttt{gfun} \\
\(E_y(z)\), ExB/profile-drive coefficient & \texttt{efun} \\
\((D_x,D_y)(z)\), magnetic drift coefficients & \texttt{dfun[:,0:2]} \\
\((g^{yy},g^{xy},g^{xx})(z)\) & \texttt{little\_g[:,0:3]} \\
\(w_z\) & \texttt{ints} \\
\bottomrule
\end{tabular}
\end{center}

The perpendicular wavenumber is evaluated as
\begin{equation}
  \kperp^2(z;k_x,k_y)
  =
  k_x^2 g^{xx}(z)
  +2 k_x k_y g^{xy}(z)
  +k_y^2 g^{yy}(z).
  \label{eq:kperp}
\end{equation}
Numerically, enforce \(\kperp^2\ge 0\) up to a small tolerance.  Negative values beyond roundoff indicate an invalid metric or coordinate conversion.

\subsection{Stellarator geometry import}

The first imported-geometry interface uses \(\rho\) as the default radial
surface coordinate, interpreted as a minor-radius-like normalized label.  The
adapter may also carry metadata for \(\psi\) or \(x\), but the internal solver
contract uses the local radial coordinate \(x\) and stores the chosen source
coordinate as metadata.  Do not differentiate through changes of radial
coordinate convention.

For stellarator flux tubes, compute or load the following physical quantities along a field line:
\begin{equation}
\begin{split}
  \mathcal{G}_{\mathrm{phys}} =
  \{&
  B,\;
  \bvec\cdot\grad z,\;
  |\grad \psi|^2,\;
  |\grad \alpha|^2,\;
  \grad\psi\cdot\grad\alpha, \\
  &(\Bvec\times\grad B)\cdot\grad\psi,\;
  (\Bvec\times\grad B)\cdot\grad\alpha,\;
  (\bvec\times\bm\kappa)\cdot\grad\psi,\;
  (\bvec\times\bm\kappa)\cdot\grad\alpha
  \},
  \label{eq:stellarator-geometry}
\end{split}
\end{equation}
where \(\bm\kappa=\bvec\cdot\grad\bvec\).  The implementation should include an adapter
\begin{equation}
  \mathcal{G}_{\mathrm{phys}}
  \longmapsto
  \{B,F,G,E_y,D_x,D_y,g^{xx},g^{xy},g^{yy}\}
  \label{eq:geometry-adapter}
\end{equation}
with tests on analytic circular geometry before using DESC/Boozer equilibria.
The first supported import path is precomputed arrays on a sampled Boozer
field line.  DESC, SIMSOPT/Boozer objects, VMEC, or booz\_xform files should be
added as source adapters that produce the same precomputed physical-array
contract.

For coupling to DESC, the intended interface is array-based rather than a
refactor of DESC inside the gyrokinetic package.  DESC should solve or optimize
the MHD equilibrium and evaluate the flux-tube quantities in
\eqref{eq:stellarator-geometry} on the selected nodes \(z_i\):
\begin{equation}
  \mathcal{A}_{\mathrm{DESC}}
  \colon
  \{E_{\mathrm{DESC}},\rho,\alpha,z_i\}_{i=1}^{N_z}
  \longmapsto
  \left\{\mathcal{G}_{\mathrm{phys}}(z_i)\right\}_{i=1}^{N_z}.
  \label{eq:desc-array-adapter}
\end{equation}
The gyrokinetic solver then maps these arrays with
\eqref{eq:geometry-adapter}, builds the fixed-topology residual, and evaluates
the objective.  If DESC exposes JAX-compatible sensitivities of
\(\mathcal{G}_{\mathrm{phys}}\) with respect to equilibrium parameters \(p\),
the design gradient is assembled by the chain rule
\begin{equation}
  \frac{\partial \mathcal{J}}{\partial p}
  =
  \sum_i
  \frac{\partial \mathcal{J}}{\partial \mathcal{G}_{\mathrm{phys}}(z_i)}
  \frac{\partial \mathcal{G}_{\mathrm{phys}}(z_i)}{\partial p}.
  \label{eq:desc-chain-rule}
\end{equation}
This keeps file I/O, equilibrium solves, and any remeshing outside the
time-advance kernels while preserving differentiation through continuous
geometry arrays.

The direct DESC extraction path uses DESC's field-line grid construction in
\((\rho,\alpha,\zeta)\) coordinates.  For a chosen surface \(\rho_0\), field-line
label \(\alpha_0\), rotational transform \(\iota(\rho_0)\), and toroidal nodes
\(\zeta_i=z_i\), DESC supplies the Cartesian vector quantities
\[
  \Bvec,\quad \grad |B|,\quad \grad\psi,\quad \grad\alpha,\quad
  \bvec,\quad \bm\kappa,
  \quad B^\zeta=\Bvec\cdot\grad\zeta.
\]
The array adapter then forms
\begin{align}
  \bvec\cdot\grad z &= B^\zeta/|B|, &
  g^{\psi\psi} &= \grad\psi\cdot\grad\psi, \\
  g^{\alpha\alpha} &= \grad\alpha\cdot\grad\alpha, &
  g^{\psi\alpha} &= \grad\psi\cdot\grad\alpha, \\
  C_\psi^B &= (\Bvec\times\grad |B|)\cdot\grad\psi, &
  C_\alpha^B &= (\Bvec\times\grad |B|)\cdot\grad\alpha, \\
  C_\psi^\kappa &= (\bvec\times\bm\kappa)\cdot\grad\psi, &
  C_\alpha^\kappa &= (\bvec\times\bm\kappa)\cdot\grad\alpha.
  \label{eq:desc-array-contractions}
\end{align}
These are the physical arrays expected by \eqref{eq:stellarator-geometry}.
A CLI fixture generator in \path{scripts/} loads either a DESC example
equilibrium or a DESC HDF5/pickle path and writes the sampled arrays to an
\path{.npz} file.

The same internal geometry can be exported to the GX/GS2 \texttt{eik}-style
contract used for stellarator parity checks:
\begin{equation}
  \mathcal{E}_{\mathrm{eik}}
  =
  \{
    \theta,\ B,\ \grad_\parallel,\ gds2,\ gds21,\ gds22,
    gbdrift,\ gbdrift0,\ cvdrift,\ cvdrift0
  \}.
  \label{eq:eik-export-contract}
\end{equation}
The implemented export uses
\[
  B=bmag,\qquad
  \grad_\parallel=F,\qquad
  (gds2,gds21,gds22)=(g^{yy},g^{xy},g^{xx}),
\]
and preserves the summed drift coefficients through
\[
  gbdrift+cvdrift=D_y,\qquad
  gbdrift0+cvdrift0=D_x.
\]
The split between grad-\(B\) and curvature drift pieces is source-code
specific and is not inferred by the exporter.  The internal mirror-force
coefficient \(G\) is therefore checked separately unless an external reference
provides a compatible proxy.  This convention lets DESC-generated flux-tube
arrays pass through the same metric/drift and \(\kperp^2\) contract used by
GX/GS2 \texttt{eik} fixtures without claiming a curvature split that was not
provided.
For GIST/GS2 text fixtures from the GX VMEC test suite, the drift columns are
read in the order
\[
  \theta,\ B,\ \grad_\parallel,\ gds2,\ gds21,\ gds22,
  cvdrift,\ cvdrift0,\ gbdrift,\ gbdrift0,
\]
and the solver compares both the individual metric fields and the summed drift
fields used by the residual.

\subsection{Magnetic drift coefficient}

Use a geometry-factorized magnetic drift frequency
\begin{equation}
  \omega_{d,s}(z,\vpar,\mu,k_x,k_y)
  =
  \frac{v_\parallel^2+\mu B(z)}{Z_s}
  \left[
    k_x D_x(z)+k_y D_y(z)
  \right],
  \label{eq:omega-d}
\end{equation}
where all dimensional constants are absorbed into the normalization and the geometry coefficients.  This is the first implementation form because it matches the Gyaradax/GKW coefficient structure and is easy to test.  Later, if curvature and grad-\(B\) drifts are separated, retain \eqref{eq:omega-d} as the assembled operator interface.

\section{Maxwellian, FLR Factors, and Polarization}

The normalized background Maxwellian is
\begin{equation}
  F_{M,s}(\vpar,\mu,z)
  =
  \frac{n_s}{(\pi T_s)^{3/2}}
  \exp\!\left[
    -\frac{\vpar^2+2\mu B(z)}{T_s}
  \right],
  \label{eq:maxwellian}
\end{equation}
where \(T_s\) and \(n_s\) are normalized species values.  Define the normalized particle energy and thermodynamic drive factor
\begin{equation}
  \mathcal{E}_s(\vpar,\mu,z)
  =
  \frac{\vpar^2+2\mu B(z)}{T_s},
  \qquad
  \Xi_s
  =
  \omega_{n,s}
  +
  \omega_{T,s}\left(\mathcal{E}_s-\frac{3}{2}\right).
  \label{eq:drive-factor}
\end{equation}

The gyroaverage of \(\phi\) is algebraic in Fourier space:
\begin{equation}
  \avg{\phi}_s
  =
  \Jzero(a_s)\phi,
  \qquad
  a_s
  =
  \frac{m_s v_{\mathrm{th},s}}{|Z_s|}
  \frac{\kperp}{B}
  \sqrt{2\mu B}.
  \label{eq:j0}
\end{equation}
The polarization factor uses
\begin{equation}
  b_s
  =
  \frac{1}{2}
  \left(
    \frac{m_s v_{\mathrm{th},s}}{Z_s}
    \frac{\kperp}{B}
  \right)^2,
  \qquad
  \Gzero(b_s) = I_0(b_s)\ee^{-b_s}.
  \label{eq:gamma0}
\end{equation}
Implementation detail: use a numerically stable exponentially scaled Bessel function for \(\Gzero\), e.g. \texttt{i0e}, and test the small-\(b\) limit \(\Gzero(b)=1-b+\mathcal{O}(b^2)\).

\section{Linear Electrostatic Gyrokinetic Model}

The first implementation evolves the collisionless linear electrostatic system
\begin{equation}
  \partial_t f_s
  =
  \RHS_s(f_s,\phi)
  =
  \mathcal{T}_{I,s}
  +\mathcal{T}_{II,s}
  +\mathcal{T}_{IV,s}
  +\mathcal{T}_{V,s}
  +\mathcal{T}_{VII,s}
  +\mathcal{T}_{VIII,s}
  +\Dop_s.
  \label{eq:linear-rhs}
\end{equation}
Term numbers follow the GKW/Gyaradax convention.  Term III, the nonlinear ExB term, is omitted in the linear baseline and documented in Section~\ref{sec:nonlinear-extension}.

Let
\begin{equation}
  a_{\parallel,s}(z,\vpar)=v_{\mathrm{th},s}F(z)\vpar,
  \qquad
  a_{\mu,s}(z,\mu)=v_{\mathrm{th},s}\mu B(z)G(z).
  \label{eq:adv-speeds}
\end{equation}
Then the six active linear terms are:
\begin{align}
  \mathcal{T}_{I,s}
  &=
  -a_{\parallel,s}\,\partial_z f_s,
  &&\hbox{parallel streaming,}
  \label{eq:term-I}
  \\
  \mathcal{T}_{II,s}
  &=
  -\ii\,\omega_{d,s}\,f_s,
  &&\hbox{magnetic drift advection,}
  \label{eq:term-II}
  \\
  \mathcal{T}_{IV,s}
  &=
  a_{\mu,s}\,\partial_{\vpar} f_s,
  &&\hbox{mirror force/trapping,}
  \label{eq:term-IV}
  \\
  \mathcal{T}_{V,s}
  &=
  \ii\,k_y\,E_y(z)\,
  J_{0,s}\phi\,
  F_{M,s}\,\Xi_s,
  &&\hbox{equilibrium-gradient drive,}
  \label{eq:term-V}
  \\
  \mathcal{T}_{VII,s}
  &=
  -\frac{Z_s}{T_s}\,
  a_{\parallel,s}\,
  F_{M,s}\,
  \partial_z(J_{0,s}\phi),
  &&\hbox{parallel field drive,}
  \label{eq:term-VII}
  \\
  \mathcal{T}_{VIII,s}
  &=
  -\frac{Z_s}{T_s}\,
  \ii\,\omega_{d,s}\,
  F_{M,s}\,
  J_{0,s}\phi,
  &&\hbox{drift field drive.}
  \label{eq:term-VIII}
\end{align}
All products are pointwise over the discrete phase-space grid after broadcasting over species, velocity, field-line, and Fourier axes.

\subsection{Dissipation}

The baseline solver includes only numerical dissipation:
\begin{equation}
  \Dop_s
  =
  \Dop_{z,s}
  +\Dop_{v_\parallel,s}
  +\Dop_{\perp,s}.
  \label{eq:diss-total}
\end{equation}
For GKW/Gyaradax-style parallel recurrence control, keep the damping inside the
residual rather than applying a post-step filter:
\begin{equation}
  \Dop_{z,s} f_s
  =
  -\texttt{disp\_par}\,
  |a_{\parallel,s}^{\rm rms}(z)|\,
  \frac{\Delta z^3}{12}\,
  \partial_z^4 f_s .
  \label{eq:gkw-disp-par}
\end{equation}
On a uniform finite-difference fallback this is exactly the GKW stencil
\([ -1,4,-6,4,-1]/(12\Delta z)\); on the spectral backend it is the
low-wavenumber-equivalent fourth derivative with the same resolution scaling.
Additional modal filtering is treated only as an experimental diagnostic.  In a
coordinate \(\xi\in\{z,\vpar,\mu\}\),
\begin{equation}
  \Dop_{\xi,s} f_s
  =
  -\nu_\xi
  T_\xi^{-1}
  \diag\left[
    \left(\frac{r}{r_{\max}}\right)^{p_\xi}
  \right]
  T_\xi f_s,
  \qquad
  r=0,\ldots,r_{\max},
  \label{eq:spectral-diss}
\end{equation}
where \(T_\xi\) is the Fourier or Chebyshev modal transform.  Set
\(\nu_\xi=0\) by default.  GKW-style fourth-order damping is used only in the
finite-difference fallback backend.

The perpendicular directions remain Fourier, so perpendicular hyper-dissipation is
\begin{align}
  \Dop_{\perp,s} f_s
  &=
  -\left[
    \nu_x\left(\frac{|k_x|}{k_{x,\max}}\right)^{p_x}
    +
    \nu_y\left(\frac{|k_y|}{k_{y,\max}}\right)^{p_y}
  \right]f_s,
  \label{eq:perp-diss}
  \\
  p_x,p_y&=
  \begin{cases}
    4, & \nu_x,\nu_y \ge 0,\\
    2, & \nu_x,\nu_y < 0\quad \hbox{(GKW-style option).}
  \end{cases}
\end{align}
For the first implementation, set default \(\nu_x=\nu_y=0\) unless a benchmark specifies otherwise.  This avoids accidental damping of Rosenbluth-Hinton and linear ITG benchmarks.

\section{Quasineutrality}

\subsection{Adiabatic electrons}

For ion species \(i\) with adiabatic electrons, solve quasineutrality independently at each \((z,k_x,k_y)\):
\begin{equation}
  \phi
  =
  -
  \frac{
    \displaystyle
    \sum_i
    Z_i n_i
    \int \dd\vpar\,\dd\mu\,
    B\,J_{0,i} f_i
  }{
    \displaystyle
    \sum_i
    \frac{Z_i^2 n_i}{T_i}
    \left[\Gamma_{0,i}-1\right]
    -
    \frac{n_e}{T_e}
  }.
  \label{eq:phi-adiabatic}
\end{equation}
The discrete numerator uses the velocity weights \(w_v,w_\mu\):
\begin{equation}
  N_i(z,k_x,k_y)
  =
  Z_i n_i
  \sum_{\ell,m}
  w_{v,\ell} w_{\mu,m}
  B(z)\,
  J_{0,i}(\ell,m,z,k_x,k_y)
  f_i(\ell,m,z,k_x,k_y).
  \label{eq:phi-numerator-discrete}
\end{equation}

For the zonal mode \(k_y=0\), GKW/Gyaradax apply a flux-surface correction to the adiabatic response.  The implementation must keep this as a separately tested path.  A practical contract is:
\begin{equation}
  \phi_{k_y=0}
  =
  \Phi_{\mathrm{local}}(f)
  +
  \Phi_{\mathrm{zf\ correction}}(f),
  \label{eq:zonal-correction-contract}
\end{equation}
where \(\Phi_{\mathrm{zf\ correction}}\) is matched against \path{linear_terms.F90} routines \path{poisson_dia} and \path{poisson_zf} and against Gyaradax tests.

\subsection{Kinetic-electron extension}

For fully kinetic species, use
\begin{equation}
  \phi
  =
  -
  \frac{
    \displaystyle
    \sum_s
    Z_s n_s
    \int \dd\vpar\,\dd\mu\,
    B\,J_{0,s} f_s
  }{
    \displaystyle
    \sum_s
    \frac{Z_s^2 n_s}{T_s}
    \left[\Gamma_{0,s}-1\right]
  }.
  \label{eq:phi-kinetic}
\end{equation}
The zero-denominator zonal/constant mode must be regularized explicitly and tested.  This path is not required before the adiabatic baseline passes tests.  The present multispecies implementation passes an algebraic TEM-favorable preflight covering charge neutrality, kinetic quasineutrality, finite coupled residuals, electron thermal streaming scaling, and the explicit CFL estimate.  This is not yet TEM validation: a converged growth rate, electron-diamagnetic-direction frequency, and complex mode structure must still agree with a revision-pinned independent solver.

\section{Discrete Grids and Operators}

The target backend is spectral in all non-perpendicular coordinates.  Production kernels should therefore be written in terms of precomputed spectral quadrature weights and differentiation matrices.  A finite-difference backend may implement the same public interface for GKW parity and fallback convergence studies.

\subsection{Velocity-space spectral collocation}

The first spectral velocity backend should keep the finite velocity box used by GKW-style benchmarks, but replace uniform finite differences with Chebyshev--Lobatto collocation.  Let
\begin{equation}
  \xi_\ell = -\cos\frac{\pi\ell}{N_v-1},
  \qquad
  v_{\parallel,\ell}=v_{\max}\xi_\ell,
  \qquad
  \ell=0,\ldots,N_v-1,
  \label{eq:vpar-cheb-grid}
\end{equation}
and
\begin{equation}
  \eta_m=\frac{1-\cos\frac{\pi m}{N_\mu-1}}{2},
  \qquad
  \mu_m=\mu_{\max}\eta_m,
  \qquad
  m=0,\ldots,N_\mu-1.
  \label{eq:mu-cheb-grid}
\end{equation}
Let \(W^C_N\) and \(D^C_N\) denote Clenshaw--Curtis quadrature weights and the Chebyshev--Lobatto first-derivative matrix on \([-1,1]\).  The physical weights and derivative matrices are
\begin{equation}
  w_{v,\ell}=v_{\max} W^C_{N_v,\ell},
  \qquad
  w_{\mu,m}=\mu_{\max} W^{C,[0,1]}_{N_\mu,m},
  \qquad
  D_{v_\parallel}=\frac{1}{v_{\max}}D^C_{N_v},
  \qquad
  D_\mu=\frac{1}{\mu_{\max}}D^{C,[0,1]}_{N_\mu}.
  \label{eq:velocity-spectral-ops}
\end{equation}
Here \(W^{C,[0,1]}\) and \(D^{C,[0,1]}\) are the mapped weights and derivative matrix on \([0,1]\).  Velocity integrals use
\begin{equation}
  \int_{-v_{\max}}^{v_{\max}}\dd\vpar
  \int_0^{\mu_{\max}}\dd\mu\,
  B(z)\,g(\vpar,\mu,z)
  \approx
  B(z)
  \sum_{\ell,m} w_{v,\ell}w_{\mu,m}
  g_{\ell m}(z).
  \label{eq:velocity-spectral-quadrature}
\end{equation}
The gyrophase factor is absorbed into the chosen normalization; if a benchmark requires an explicit \(2\pi\), it should enter the documented velocity measure, not an individual physics term.  Velocity collocation is the selected production CPU representation: Chebyshev is the general reduced/differentiable path, GKW finite differences retain convention parity, and the midpoint/Gauss--Laguerre recipe is used by the source-matched W7-X gate.  Hermite--Laguerre remains a reduced GX discriminator until it is integrated with the complete stellarator residual and field solve and passes comparable convergence, parity, gradient, and timing gates.

\subsection{Parallel spectral grid}

The parallel coordinate \(z\) is the chosen field-line parameter.  For a closed mode chain, use Fourier collocation along the concatenated connected coordinate:
\begin{equation}
  z_j=z_0+\frac{jL_z}{N_z},
  \qquad
  q_r=\frac{2\pi}{L_z}
  \begin{cases}
    r, & 0\le r\le N_z/2,\\
    r-N_z, & N_z/2<r<N_z,
  \end{cases}
  \label{eq:z-fourier-grid}
\end{equation}
and
\begin{equation}
  D_z = F_z^{-1}\diag(\ii q_r)F_z,
  \qquad
  w_{z,j}=L_z/N_z.
  \label{eq:z-fourier-derivative}
\end{equation}
Here \(F_z\) is the discrete Fourier transform along the connected field-line chain.  For an open or non-closing chain, use Chebyshev--Lobatto collocation on the finite interval and the corresponding derivative matrix.  In both cases the streaming operator is \(a_{\parallel,s}D_z f_s\); the geometry object supplies \(\bvec\cdot\grad z\), metric factors, and quadrature weights for the chosen coordinate.

For imported stellarator geometry, interpolate or evaluate \(B\), metric terms, and drift coefficients on the same spectral \(z\) nodes.  Do not hide the quadrature rule in the RHS; store \(z_j\), \(w_{z,j}\), \(D_z\), and any modal transform needed by filters in the geometry/precompute object.

\subsection{Fourier grids}

The \(k_y\) grid is a nonnegative half spectrum:
\begin{equation}
  k_{y,n} = n\Delta k_y,\qquad n=0,\ldots,N_y-1.
  \label{eq:ky-grid}
\end{equation}
The \(k_x\) grid is centered:
\begin{equation}
  k_{x,m} = (m-m_0)\Delta k_x,\qquad m=0,\ldots,N_x-1,
  \label{eq:kx-grid}
\end{equation}
where \(m_0\) indexes \(k_x=0\).  In single-mode linear runs, \(N_x=1\) and \(k_x=0\) is allowed.

Parseval weights for the \(k_y\) half spectrum are
\begin{equation}
  P(k_y)=
  \begin{cases}
    1, & k_y=0,\\
    2, & k_y>0.
  \end{cases}
  \label{eq:parseval}
\end{equation}

\subsection{Twist-and-shift connectivity}

At a parallel boundary, nonzonal modes connect to a shifted radial wavenumber.  Represent this with integer maps
\begin{equation}
  \mathrm{ixplus}(m,n),\qquad
  \mathrm{ixminus}(m,n),
  \label{eq:ixmaps}
\end{equation}
where invalid entries denote an open end of a connected mode chain.  For axisymmetric tests, the shift should reproduce GKW mode-box behavior.  For stellarators, compute the shift from the field-line twist:
\begin{equation}
  k_x^{+}
  =
  k_x + k_y\,\Delta_{\mathrm{twist}},
  \qquad
  \Delta_{\mathrm{twist}}
  =
  \left.
  \frac{\partial \alpha_{\mathrm{end}}}{\partial x}
  \right|_{\psi_0}
  -
  \left.
  \frac{\partial \alpha_{\mathrm{start}}}{\partial x}
  \right|_{\psi_0},
  \label{eq:twist-shift-general}
\end{equation}
then choose \(\Delta k_x\) so the shifted mode lands on an integer grid index when possible.  The continuous twist can be differentiable, but the integer connectivity maps are static topology and must not be differentiated.

\subsection{Spectral derivative contract}

All derivative-using RHS terms should call backend operators:
\begin{equation}
  \partial_z f \mapsto D_z f,
  \qquad
  \partial_{\vpar} f \mapsto D_{v_\parallel} f,
  \qquad
  \partial_\mu f \mapsto D_\mu f.
  \label{eq:spectral-derivative-contract}
\end{equation}
The implementation may store nodal-to-modal transforms \(T_\xi\) for filtering and diagnostics:
\begin{equation}
  \widehat f^{(\xi)} = T_\xi f,
  \qquad
  f = T_\xi^{-1}\widehat f^{(\xi)}.
  \label{eq:modal-transform}
\end{equation}
For smooth manufactured functions, derivative tests should demonstrate spectral convergence until floating-point or interpolation error dominates.  On nonsmooth or open-boundary tests, the code should still converge monotonically and the finite-difference backend provides a useful sanity check.

\subsection{GKW finite-difference fallback}

The following GKW/Gyaradax-style formulas are fallback operators, not the primary target discretization from \path{task.tex}.  They are useful for parity tests against the original implementation and for debugging spectral boundary behavior.

Interior fourth-order centered first derivative:
\begin{equation}
  (\partial_\xi f)_j
  =
  \frac{
    f_{j-2}-8f_{j-1}+8f_{j+1}-f_{j+2}
  }{12\Delta \xi}
  +\mathcal{O}(\Delta\xi^4).
  \label{eq:d1-centered}
\end{equation}
Interior fourth derivative stencil:
\begin{equation}
  (\partial_\xi^4 f)_j
  =
  \frac{
    f_{j-2}-4f_{j-1}+6f_j-4f_{j+1}+f_{j+2}
  }{\Delta\xi^4}
  +\mathcal{O}(\Delta\xi^2).
  \label{eq:d4-centered}
\end{equation}

The fallback parallel streaming operator uses sign-dependent upwinding.  Boundary stencils follow GKW/Gyaradax.  The implementation should store dimensionless stencil tables
\begin{equation}
  S_{z,+}^{(1)},\quad S_{z,-}^{(1)},\quad
  S_{z,+}^{(4)},\quad S_{z,-}^{(4)}
  \label{eq:stencil-tables}
\end{equation}
and select them by \(\operatorname{sign}(a_{\parallel,s})\).  At a boundary, a stencil point outside the \(z\)-domain is resolved with \(\mathrm{ixplus}\) or \(\mathrm{ixminus}\); if no connected mode exists, the point is treated with the chosen open-boundary rule.

For the fallback mirror term in \(v_\parallel\), use the centered fourth-order derivative in the interior and zero-padding/no-flux style boundary behavior, matching GKW velocity-space conventions.  Every derivative routine must have manufactured-solution tests in both spectral and fallback modes.

For production GKW parity with \(\texttt{ltrapping\_arakawa}\), the fallback
also exposes a fused matrix-free \(\texttt{igh}\) operator.  This operator does
not apply Term I and Term IV as two independent derivatives.  Instead it
precomputes row-local source-shift weights \(C_{\Delta v,\Delta z}\) with
\(\Delta v,\Delta z\in[-2,2]\) and applies
\begin{equation}
  \left(\mathcal{L}_{\mathrm{igh}} f\right)_{a,b,j,k_x,k_y}
  =
  \sum_{\Delta v=-2}^{2}
  \sum_{\Delta z=-2}^{2}
  C_{\Delta v,\Delta z;a,b,j,k_x,k_y}\,
  f_{a+\Delta v,b,j+\Delta z,k_x',k_y}.
  \label{eq:gkw-igh-backend}
\end{equation}
The shifted \(k_x'\) is obtained from the same GKW mode-connectivity maps used
by the upwind fallback.  The coefficients include the Arakawa Hamiltonian
stencil for \(H=v_\parallel^2/2+\mu B\) and the in-operator
\(\texttt{disp\_par}\) and \(\texttt{disp\_vp}\) fourth-difference diffusion
terms.  This backend is a parity/debug path for the GKW finite-difference
grid; the differentiable spectral target remains the default solver route.

\section{Residual Assembly and Precomputation}

\subsection{Precomputed coefficient arrays}

Before a time loop or eigenvalue solve, construct a precompute object
\begin{equation}
  \mathcal{P} =
  \left\{
  \begin{array}{l}
  J_{0,s},\Gamma_{0,s},F_{M,s},\Xi_s,\omega_{d,s},\\
  a_{\parallel,s},a_{\mu,s},\\
  D_z,D_{v_\parallel},D_\mu,\quad
  T_z,T_{v_\parallel},T_\mu,\\
  \hbox{fallback stencil tables},\\
  \hbox{phi weights},\quad
  \hbox{diagnostic weights}
  \end{array}
  \right\}.
  \label{eq:precompute}
\end{equation}
Shapes for the adiabatic single-species path should be
\begin{equation}
\begin{array}{ll}
f: & (N_v,N_\mu,N_z,N_x,N_y),\\
\phi: & (N_z,N_x,N_y),\\
\Jzero,F_M,\omega_d: & (N_v,N_\mu,N_z,N_x,N_y),\\
a_\parallel,a_\mu: & (N_v,N_\mu,N_z,1,1).
\end{array}
\end{equation}
Multi-species arrays add a leading species axis.

\subsection{Public solver interfaces}

The core interfaces should be pure functions:
\begin{align}
  \phi &= \texttt{calculate\_phi}(f,\mathcal{G},p,\mathcal{P}),
  \label{eq:api-phi}
  \\
  r &= \texttt{linear\_residual}(f,\phi,\mathcal{G},p,\mathcal{P}),
  \label{eq:api-residual}
  \\
  f^{n+1} &= \texttt{rk4\_step}(f^n,\mathcal{G},p,\mathcal{P}).
  \label{eq:api-rk4}
\end{align}
No file I/O, Python-side mutation, random initialization, or diagnostic printing belongs inside these functions.  This keeps the residual JIT-compatible and differentiable.

\section{Time Integration and Linear Growth Rates}

\subsection{RK4}

The baseline time integrator is explicit fourth-order Runge--Kutta:
\begin{align}
  k_1 &= \RHS(f^n),\\
  k_2 &= \RHS(f^n+\tfrac{1}{2}\Delta t\,k_1),\\
  k_3 &= \RHS(f^n+\tfrac{1}{2}\Delta t\,k_2),\\
  k_4 &= \RHS(f^n+\Delta t\,k_3),\\
  f^{n+1} &= f^n + \frac{\Delta t}{6}(k_1+2k_2+2k_3+k_4).
  \label{eq:rk4}
\end{align}
Each RHS evaluation recomputes \(\phi=\Phi(f)\).  Multi-step runs should use a scan-style loop so that JAX can compile the full step block.

\subsection{Timestep estimate}

Use fixed \(\Delta t\) first.  Add adaptive estimates after the fixed-step solver passes benchmarks.  For the spectral backend, estimate the explicit stability limit using operator radii:
\begin{equation}
  \rho_1 =
  \rho\!\left(a_\parallel D_z+a_\mu D_\mu\right),
  \qquad
  \rho_{\mathrm{diss}} =
  \rho\!\left(\Dop_z+\Dop_{v_\parallel}+\Dop_\mu+\Dop_\perp\right),
  \label{eq:spectral-radii}
\end{equation}
where \(\rho(\cdot)\) is either an inexpensive power-iteration estimate or an upper bound from row sums on small grids.  A practical RK4 estimate is
\begin{equation}
  \Delta t_{\mathrm{spec}}
  =
  c_{\mathrm{safety}}
  \min\left(
    \frac{2.4}{\max(\rho_1,\epsilon)},
    \frac{2.4}{\max(\rho_{\mathrm{diss}},\epsilon)}
  \right).
  \label{eq:dtspec}
\end{equation}

The finite-difference fallback may also expose the GKW-compatible linear estimate
\begin{equation}
  \Delta t_{\mathrm{lin}}
  =
  \frac{c_{\mathrm{safety}}}
  {
    \max\left[
      40,\;
      \frac{t_1}{2.4},\;
      \frac{t_4}{2.4},\;
      \frac{t_{\mathrm{field}}}{2.4}
    \right]
  },
  \label{eq:dtlin}
\end{equation}
with
\begin{align}
  t_1 &=
  \max\left(
    c_{z,1}\frac{\max|a_\parallel|}{\Delta z},
    c_{v,1}\frac{\max|a_\mu|}{\Delta v}
  \right),
  \label{eq:t1}
  \\
  t_4 &=
  \max\left(
    c_{z,4}\nu_z\frac{\max|a_\parallel|}{\Delta z},
    c_{v,4}\nu_v\frac{\max|a_\mu|}{\Delta v}
  \right).
  \label{eq:t4}
\end{align}
Here \(2.4\) is the RK4 stability-boundary factor used in GKW/Gyaradax notes.  \(t_{\mathrm{field}}\) is needed mainly for kinetic electrons and can be zero in the adiabatic baseline.  The effective timestep is
\begin{equation}
  \Delta t
  =
  \min(\Delta t_{\mathrm{input}},\Delta t_{\mathrm{spec}},\Delta t_{\mathrm{lin}},\Delta t_{\mathrm{NL}}),
  \label{eq:dteff}
\end{equation}
where \(\Delta t_{\mathrm{lin}}\) is ignored unless the fallback backend is active, and \(\Delta t_{\mathrm{NL}}\) is relevant only after the nonlinear ExB term is implemented.

\subsection{Mode amplitude, growth rate, and frequency}

Let \(\mathcal{C}(k_y)\) be the connected \(k_x\)-chain containing \(k_x=0\).  Define the per-\(k_y\) potential amplitude
\begin{equation}
  A(k_y,t)
  =
  \left[
    \sum_{z} w_z
    \sum_{k_x\in\mathcal{C}(k_y)}
    |\phi(z,k_x,k_y,t)|^2
  \right]^{1/2}.
  \label{eq:mode-amplitude}
\end{equation}
Over a diagnostic window \([t_a,t_b]\),
\begin{equation}
  \gamma(k_y)
  =
  \frac{\log A(k_y,t_b)-\log A(k_y,t_a)}{t_b-t_a}.
  \label{eq:growth-rate}
\end{equation}
A robust real-frequency estimator is the phase rotation of the complex mode:
\begin{equation}
  \omega(k_y)
  =
  -\frac{1}{t_b-t_a}
  \arg
  \left[
    \sum_z w_z
    \sum_{k_x\in\mathcal{C}(k_y)}
    \phi^*(z,k_x,k_y,t_a)\phi(z,k_x,k_y,t_b)
  \right].
  \label{eq:frequency}
\end{equation}
In linear initial-value runs, normalize each active \(k_y\) after each diagnostic window:
\begin{equation}
  f_s(\cdots,k_y) \leftarrow
  \frac{f_s(\cdots,k_y)}{\max(A(k_y),\epsilon_A)}.
  \label{eq:linear-normalization}
\end{equation}
Store the accumulated normalization factor so diagnostics remain physically interpretable.

\section{Fluxes and Quasilinear Objective}

The immediate optimization objective is not nonlinear heat flux, but a differentiable linear proxy.  For each \(k_y\), define a mode-structure weighted perpendicular scale
\begin{equation}
  \avg{\kperp^2}_{k_y}
  =
  \frac{
    \displaystyle
    \sum_z w_z
    \sum_{k_x\in\mathcal{C}(k_y)}
    \kperp^2(z,k_x,k_y)|\phi(z,k_x,k_y)|^2
  }{
    \displaystyle
    \sum_z w_z
    \sum_{k_x\in\mathcal{C}(k_y)}
    |\phi(z,k_x,k_y)|^2+\epsilon_A
  }.
  \label{eq:kperp-average}
\end{equation}
The first quasilinear proxy is
\begin{equation}
  Q_{\mathrm{QL}}
  =
  \sum_{k_y>0}
  W(k_y)
  \frac{\max(\gamma(k_y),0)}{\avg{\kperp^2}_{k_y}+\epsilon_k}.
  \label{eq:ql-proxy}
\end{equation}
For smooth reverse-mode differentiation, replace \(\max(\gamma,0)\) by a softplus if needed:
\begin{equation}
  \max(\gamma,0)
  \approx
  \tau\log\left(1+\ee^{\gamma/\tau}\right).
  \label{eq:softplus-growth}
\end{equation}
This objective matches the spirit of direct microstability optimization while remaining cheap enough for early design-loop tests.

\section{Differentiability Contract}

The following quantities should be differentiable:
\begin{itemize}
  \item continuous geometry arrays \(B,F,G,E_y,D_x,D_y,g^{xx},g^{xy},g^{yy}\);
  \item profile gradients \(\omega_n,\omega_T\);
  \item species parameters where physically meaningful;
  \item scalar geometry knobs in analytic tests, e.g. \(q,\hat s,\epsilon\);
  \item objective functions \(A,\gamma,\omega,Q_{\mathrm{QL}}\) through a fixed number of time steps.
\end{itemize}
The following quantities are static/non-differentiable:
\begin{itemize}
  \item grid sizes;
  \item integer mode labels and connectivity maps;
  \item branch choices for topology;
  \item file I/O and equilibrium loading;
  \item random initialization keys.
\end{itemize}
For stellarator optimization, a differentiable equilibrium code may provide geometry arrays and their sensitivities.  The gyrokinetic solver should not require differentiating through an integer remeshing or topology change.
The implementation enforces this boundary with a versioned topology contract
that fingerprints velocity and parallel nodes and backends, perpendicular mode
sets, twist-and-shift connectivity, and provider field-line metadata.  A changed
fingerprint invalidates compiled objectives, checkpoints, and gradients.
In the present implementation, imported DESC-style arrays are accepted as a
\texttt{desc} geometry model in the single-surface objective; the direct DESC object
adapter is optional and remains outside the solver-core dependency chain.


\section{Current Implementation and Validation Status}
\label{sec:current-status}

The implementation now covers the linear electrostatic architecture specified
above.  The core package in \path{src/stellarator_gk/} contains PyTree parameter
objects, Fourier grids, spectral and finite-difference derivative backends,
analytic and imported stellarator geometry adapters, quasineutrality solvers,
term-level linear RHS assembly, RK4 time advance, growth/frequency diagnostics,
stable design objectives, gradient audits, robust multi-sample aggregation,
fixed-topology guards, and schema-versioned optimization checkpoints.  The production
physics target remains linear electrostatic gyrokinetics with adiabatic
electrons; nonlinear turbulence and extended physics are deferred to
Sec.~\ref{sec:nonlinear-extension}.

\subsection{Validation guardrails}

The current regression suite is intentionally layered.  Algebraic tests cover
broadcasting, PyTree behavior, derivative operators, field solves, term
linearity, and differentiability on reduced grids.  Code-to-code and benchmark
fixtures then check the pieces most likely to hide convention errors:
\begin{itemize}
  \item the true Rosenbluth--Hinton late-plateau metric passes against the
  GKW/Gyaradax value \(0.0711\), with observed value \(0.070413\);
  \item the Cyclone selected-\(k_y\) scalar growth gate passes after matching
  GKW's internal wave-number convention;
  \item GKW selected-mode RHS/action traces, imported-state one-window replays,
  initial-state normalization, and row-normalized \path{parallel_phi.dat}
  profile checks pass at their current tolerances;
  \item eik geometry checks pass for DESC/GX fixtures, solver export, GX/GS2
  import, and independent GX/VMEC GIST sources;
  \item reduced stellarator scan and DESC optimization examples execute through
  fixed-topology differentiable objectives, but are not production claims.
\end{itemize}

These passing gates are guardrails, not the end state.  Two confidence gaps are
kept visible: full selected-mode GKW state-history parity remains open, and the
multi-time GKW velocity-slice error grows from \(3.99\times10^{-3}\) at step 20
to \(3.67\times10^{-2}\) by step 800.  They do not currently block the W7-X
external-reference path, but they remain useful warnings against overclaiming.

\subsection{One-command stellarator scans}

The reduced imported-geometry scan is executable with
\begin{verbatim}
uv run python examples/run_stellarator_linear_scan.py
\end{verbatim}
By default it loads \path{fixtures/desc_geometry_dshape_rho05_alpha0.npz}, runs
a short fixed-topology linear scan, and writes geometry audits, growth/frequency
tables, mode structures, convergence metadata, and a quasilinear proxy.  The
same command can read DESC paths, eik tables, and stella \path{.geometry}
files.  Before solving, it applies the shared geometry-preflight contract:
finite fields, positive \(B\), positive metric diagonals, representative
nonnegative \(k_\perp^2\), GX/GS2 eik-export compatibility, and an optional
finite-difference mirror-force check when the source data supports it.  The
separate mode-boundary contract checks zonal identity maps, nonzonal open chain
ends, reciprocal linked shifts, and \texttt{ikxspace} spacing.

A real-geometry reduced W7-X fixture is stored in
\path{fixtures/w7x_itg_reduced_benchmark/}.  It uses the local GX/GIST W7-X eik
geometry and GX W7-X adiabatic-electron ITG input as provenance.  Its scalar
outputs are finite regression artifacts, not an external growth-rate claim; the
fixture is deliberately guarded by production-readiness ledgers.

\subsection{Matched W7-X stella reference}

The preferred external W7-X reference is stella, a CPU continuum gyrokinetic
code that can run the matched case locally.  Acceptance uses
\(\texttt{torflux}=0.64\), \(\alpha=0\), adiabatic electrons, \(k_x=0\),
\(k_y=0.3\), \(N_z=256\), \(N_{v_\parallel}=32\), \(N_\mu=8\),
\(\Delta t=0.1\), and \(t=500\).  Generated stella input, geometry, state, and
output remain external scratch data rather than canonical repository assets.
The lower \(k_y=0.1,0.2\) branches are diagnostic because their stella omega
windows fail the convergence test even at \(t=1000\).

The solver imports the external stella \path{.geometry} table only for this
independent validation path.  With the source-matched analytic initial
perturbation, the converged branch gives
\[
(|\Delta\gamma|,|\Delta\omega_r|,\epsilon_\phi)
=(3.20\times10^{-4},2.18\times10^{-3},1.08\times10^{-2}),
\]
where \(\epsilon_\phi\) is the phase-aligned complex profile error.  All three
values pass the \(2\times10^{-2}\) gate.

\subsection{Resolved W7-X timestep discriminator}

Source instrumentation established the production ordering, which differs from
a generic RK4 Strang split.  Each stella step consists of three SSP RK3 explicit
stages, one full sign-dependent cubic semi-Lagrangian mirror characteristic,
and one full implicit parallel-streaming and quasineutrality response.  On an
exact native stella input, the JAX replay errors after RK3, after mirror advance,
in the final distribution, and in potential are respectively
\[
7.19\times10^{-6},\quad 1.34\times10^{-4},\quad
2.17\times10^{-5},\quad 1.34\times10^{-4}.
\]
This closes the timestep discriminator independently of long-time eigenmode
selection.

The earlier long-time profile mismatch came from comparing different initial
mode mixtures.  Stella initializes a localized Maxwellian perturbation,
\[
g_0=A(1+i)\exp[-(2\pi z)^2]
\exp[-v_\parallel^2-2\mu B(z)],
\]
whereas the generic scan used a dense oscillatory state.  The
\texttt{stella\_maxwellian} option implements this expression analytically; no
W7-X distribution is stored.  This distinction matters at finite time for a
nonnormal spectrum with slowly separating branches.

\subsection{Optimization status}

The reduced fixed-topology optimization milestone is complete.  The public
objective contract composes selected, hard-maximum, or smooth-maximum growth,
real-frequency targets, phase-invariant complex mode-structure penalties, and
quasilinear proxies.  Reverse-mode derivatives are checked against central
differences through geometry arrays and time advance, with explicit diagnostics
for near-degenerate \(k_y\) branches.  Fixed scans over \(\rho\), \(\alpha\), and
\(k_y\) support weighted-mean, hard-worst-case, and differentiable
soft-worst-case reduction.  Checkpoints retain the complete sample tensor,
objective policy, topology fingerprints, provider provenance, revisions,
command, seed, parameters, and iteration history.

A reduced W7-X loop loads the installed VMEC++ named configuration, performs
fresh center/plus/minus equilibrium solves for a boundary harmonic, converts
each result in memory, verifies fixed topology, and takes an outer finite-
difference step.  No equilibrium artifact is stored in the repository.  This
demonstrates the real MHD provider boundary but does not differentiate through
VMEC++, claim unrestricted equilibrium-shape optimization, or predict nonlinear
turbulent heat flux.

\subsection{Production criteria}

The trusted linear W7-X branch now satisfies:
\begin{enumerate}
  \item weighted stella-vs-solver \(k_y=0.3\) term and stage-array parity;
  \item a passing converged W7-X growth, frequency, and mode-profile gate;
  \item source-matched field-line, 256-point parallel, 32-by-8 velocity,
  timestep, split-ordering, initial-state, and growth-window controls;
  \item guarded CPU timing of the validated advance after parity passes.
\end{enumerate}
Production full-shape design optimization additionally requires a
differentiable equilibrium provider (or a validated sensitivity interface) and
physics beyond the present reduced linear objective.  A broader low-\(k_y\)
spectral claim additionally requires converged independent reference windows.


\section{Nonlinear and Extended Physics Roadmap}
\label{sec:nonlinear-extension}

The nonlinear ExB term retained for later is
\begin{equation}
  \mathcal{T}_{III,s}
  =
  -\left\{J_{0,s}\phi,f_s\right\},
  \qquad
  \{a,b\}
  =
  \partial_x a\,\partial_y b-\partial_y a\,\partial_x b.
  \label{eq:nonlinear-bracket}
\end{equation}
The implementation should be pseudospectral:
\begin{enumerate}
  \item multiply by \(\ii k_x\) and \(\ii k_y\) in Fourier space;
  \item inverse FFT to a dealiased real grid;
  \item compute the bracket in real space;
  \item FFT back and truncate to the retained spectral grid.
\end{enumerate}
The implemented nonlinear foundation follows this construction for centered
$k_x$ modes and the nonnegative-$k_y$ Hermitian half spectrum.  It also
evaluates an advective timestep bound from the padded-grid ExB velocity.  The
host-controlled nonlinear driver chooses the minimum of this instantaneous
bound and the coupled linear bound for each accepted RK4 step.  Since timestep
accept decisions are nonsmooth, differentiable trajectories retain the
fixed-step driver. Production nonlinear readiness additionally requires
statistically stationary heat flux, resolution and box-size convergence, and
independent-code parity.

For a collocation-space state, the implemented electrostatic heat response is
egin{equation}
  \mathcal{Q}_{s,\mathbf{k}}(z)
  = \int d^3v\,J_{0,s}\,T_s
    \left(\mathcal{E}_s-\frac{3}{2}\right)f_{s,\mathbf{k}},
end{equation}
and its radial spectral contribution is proportional to
(k_y\operatorname{Im}[\phi_{\mathbf{k}}^*\mathcal{Q}_{s,\mathbf{k}}]).
The velocity integral uses the same discrete measure as the field solve.  This
defines the evolved-state diagnostic but does not replace the stationarity,
resolution, domain-size, and independent-code gates.

Other extensions should be added only after the linear electrostatic baseline is validated:
\begin{itemize}
  \item kinetic electrons via \eqref{eq:phi-kinetic};
  \item collisions from \path{collisionop.F90};
  \item electromagnetic \(A_\parallel\) and \(B_\parallel\);
  \item toroidal rotation and ExB shear from \path{rotation.F90};
  \item nonlinear heat-flux objectives.
\end{itemize}

As an intermediate collision model, the collocation solver may use a
species-local linearized BGK operator
\begin{equation}
  C_s[f_s] = -\nu_s\left(f_s-\mathcal{P}_s f_s\right),
\end{equation}
where the discrete projection $\mathcal{P}_s$ spans $F_{M,s}$,
$v_\parallel F_{M,s}$, and $\mathcal{E}_s F_{M,s}$ in the same quadrature used
by the field solve.  This construction conserves the discrete density,
parallel-momentum, and energy moments species by species.  It is a
differentiable model-collision foundation, not a replacement for a validated
Landau/Fokker--Planck operator with inter-species exchange.

The finite-difference Fokker--Planck development path also retains the ordered
test-particle contributions (C_{ab}[f_a]).  Its experimental pairwise
completion groups (C_{ab}[f_a]) and (C_{ba}[f_b]), then applies a discrete
projection for the two density constraints and the pair's combined physical
parallel-momentum and energy constraints.  Consequently, the directed
contributions satisfy
egin{equation}
  int C_{ab}^{\mathrm{pair}}\,d^3v = 0,
  \qquad
  \sum_{s\in\{a,b\}} \int
  \begin{pmatrix}p_{\parallel,s}\\ E_s\end{pmatrix}
  C_{s\bar{s}}^{\mathrm{pair}}\,d^3v = 0.
end{equation}
This creates reciprocal cross-species data flow and a pairwise conservation
contract.  It is not the Laguerre--Legendre field-particle response used by
stella and must not be described as a validated Landau operator until its
coefficients and action pass an independent executable parity gate.

The validated software form for a future Laguerre--Legendre completion is the
low-rank pair/component contraction
\begin{equation}
  \Psi_{ab,c}(z,\mathbf{k})
  = \sum_{v_\parallel,\mu}
    D_{ab,c}(v_\parallel,\mu,z,\mathbf{k}) f_b,
  \qquad
  C^{\mathrm{fp}}_a
  = \sum_{b,c} R_{ab,c}\Psi_{ab,c},
\end{equation}
where $c=(l,m,j)$, $D$ includes the discrete velocity measure and all driver
normalizations, and $R$ is the target-space response.  The JAX implementation
accepts these tensors without depending on stella.  Native stella traces verify
the response application and indexing at roundoff.  The response $R$ is now
constructed independently from the physical velocity and magnetic grids,
species mass ratios, pair collision frequencies, gyroaverages, and
$\Delta_j$.  A 44,928-row native primitive trace agrees with this construction
to $1.84\times10^{-14}$ in relative $L_2$, and the analytic incomplete-gamma
construction of $\Delta_0$ agrees to $4.25\times10^{-13}$.  Independent
construction of the higher-$j$ recurrence now uses the complete native
velocity-integration measure and agrees for all traced $j=0,1$ values to
$7.84\times10^{-13}$ in scaled $L_2$.  Construction of the normalized driver
$D$ now uses the reversed-pair recurrence, Maxwellian, background gyroaverage,
velocity measure, and locally evaluated pair normalization; it agrees with
44,928 native coefficients to $3.07\times10^{-11}$ in scaled $L_2$.  The local
$D$ and $R$ tensors can therefore be assembled into the differentiable
low-rank precompute without stella.  A second state with changed initialization
retains byte-identical coefficient traces and replays the native solved-$\Psi$
action at $1.39\times10^{-13}$ relative $L_2$.  Thus the remaining distinction
is the differential test-particle matrix.  A differentiable Woodbury solve now
combines any supplied $I-\Delta t C_{\mathrm{tp}}$ with $D$ and $R$ in
stella's implicit ordering and agrees with a directly materialized dense
backward-Euler system at roundoff.  Native construction parity for
$C_{\mathrm{tp}}$ remains required.
The independently Gyaradax-matched nine-point differential stencil can now be
materialized as $I-\Delta t C_{\mathrm{tp}}$ on each spatial/Fourier batch and
inserted directly into this solve; its matrix application matches the original
stencil action at roundoff.  A stricter same-executable stella matrix gate and
production residual wiring remain before claiming full stella parity.


\section{Testing and Benchmark Contract}

Every public function in \path{src/} should have focused tests.  For a narrow
change, update the smallest relevant unit or fixture test; for a change that
alters shared physics, geometry, time advancement, or validation contracts, run
the corresponding focused gate plus the broader suite.

\subsection{Required unit-test coverage}

\begin{longtable}{p{0.28\textwidth}p{0.62\textwidth}}
\toprule
Component & Required tests \\
\midrule
Parameter PyTrees & flatten/unflatten, static/dynamic fields, JIT behavior. \\
Grids and topology & shapes, centering, quadrature weights, Parseval weights, mode labels, linked/open boundary maps. \\
Derivative backends & manufactured derivative checks, finite-difference parity, filter behavior. \\
Geometry & finite fields, \(B>0\), \(k_\perp^2\ge0\), eik/DESC/stella adapter contracts, differentiability of continuous arrays. \\
Physics primitives & Bessel/FLR limits, Maxwellian and drive factors, mirror/drift/streaming coefficients, broadcasting, AD checks. \\
Field solve & zero-input behavior, zonal and nonzonal paths, quasineutrality residuals. \\
RHS terms & zero-input behavior, linearity, isolated term formulas, full residual reconstruction, JIT compatibility. \\
Time advance & RK4 order, fixed-step scans, growth/frequency recovery, phase-unwrapped frequency on long windows. \\
Objectives & finite values, gradients, finite-difference checks on reduced fixed-topology problems. \\
Validation scripts & committed fixture schemas, status JSON contracts, failure states, and command-line smoke paths. \\
\bottomrule
\end{longtable}

\subsection{Benchmark ladder}

The benchmark order should remain:
\begin{enumerate}
  \item manufactured derivatives, field solves, and scalar ODE tests;
  \item Gyaradax circular and s-alpha parity;
  \item GKW simple/Cyclone fixtures and selected-mode RHS/action traces;
  \item Rosenbluth-Hinton late plateau;
  \item Cyclone selected-\(k_y\) scalar and mode-structure diagnostics;
  \item DESC/GX/GS2/stella eik and geometry contracts;
  \item matched W7-X stella growth, frequency, and complex \(\phi(z)\) parity;
  \item production-control W7-X convergence and CPU timing;
  \item only then DESC optimization campaigns and nonlinear extensions.
\end{enumerate}

\subsection{First trusted solver milestone}

The first trusted solver milestone is complete when \(\Phi(f)\), \(\RHS(f)\),
RK4 stepping, \(\gamma\), \(\omega\), and \(Q_{\mathrm{QL}}\) are implemented,
the reduced RH/Cyclone/GKW/DESC/eik guardrails pass, a matched W7-X stella
mode-structure gate passes, production-control convergence and CPU timing are
recorded, and \path{STATUS.md} states the remaining caveats without using
reduced examples as production evidence.

\begin{thebibliography}{9}

\bibitem{peeters2009gkw}
A.~G. Peeters, Y. Camenen, F.~J. Casson, W.~A. Hornsby, A.~P. Snodin, D. Strintzi, and G. Szepesi.
\newblock The nonlinear gyro-kinetic flux tube code GKW.
\newblock \emph{Computer Physics Communications}, 180:2650--2672, 2009.

\bibitem{gyaradax}
G. Galletti, E. Volkmann, and J. Brandstetter.
\newblock \texttt{gyaradax}: local gyrokinetics JAX code.
\newblock Project materials in \path{papers/gyaradax-paper/} and \path{relevant-codes/gyaradax/}.

\bibitem{jorge2024}
R. Jorge, W. Dorland, P. Kim, M. Landreman, N.~R. Mandell, G. Merlo, and T. Qian.
\newblock Direct microstability optimization of stellarator devices.
\newblock \emph{Physical Review E}, 110:035201, 2024.

\bibitem{kim2024}
P. Kim et al.
\newblock Optimization of nonlinear turbulence in stellarators.
\newblock \emph{Journal of Plasma Physics}, 90:905900210, 2024.

\end{thebibliography}

\end{document}
